% Trigonometrie · Seite berechnen · Prüfungs-Fokus MSA – bank/trigonometrie/e1.jsonl (zusammenbau v0.6, Prüfungs-Fokus)
\einheitenkopf[e1]{Einheit 1 · Seite berechnen mit Sinus, Kosinus und Tangens}
\zweigzeile{ab Kl. 10 (am Gymnasium ab Kl. 9) · P10 ×11}
\setcounter{aufgabe}{0}

% Anlauf (ohne Prüfkennung)
\begin{aufgabe}{Seite berechnen – Anlauf}
\begin{teile}
\teil Rechter Winkel bei $B$. Schreibe an jede Seite H, G oder A – vom Winkel $\alpha$ aus.

\dreieck{(0,0)}{(5,0)}{(5,3)}{r}{t}{s}{\alpha}{}{}
\teil Rechter Winkel bei $A$. Trage den Bruch ein.

\dreieck{(0,0)}{(4,0)}{(0,3)}{x}{y}{z}{}{}{\gamma}
$\mathrm{tan}\,\gamma =$ \leerfeld
\teil Taschenrechner im Grad-Modus: $\mathrm{sin}\,23^\circ$ auf drei Dezimalen? \leerfeld
\end{teile}
\end{aufgabe}

\begin{aufgabe}{Seite berechnen – Prüfungsaufgaben}
\begin{teile}
\teil Rechter Winkel bei $C$. Trage den Bruch ein. \hfill (P17) \\

\dreieck{(0,0)}{(5,0)}{(1.8,2.4)}{d}{e}{f}{}{\beta}{}
$\mathrm{sin}\,\beta =$ \leerfeld
\teil Rechter Winkel bei $B$. Trage den Bruch ein. \hfill (P19) \\

\dreieck{(0,0)}{(5,0)}{(5,3)}{k}{m}{n}{\alpha}{}{}
$\mathrm{cos}\,\alpha =$ \leerfeld
\teil Rechter Winkel bei $A$. Welche Gleichung gilt? \hfill (P18) \\ \kreuz{$\mathrm{sin}\,\beta = \frac{p}{q}$} \\ \kreuz{$\mathrm{cos}\,\beta = \frac{q}{p}$} \\ \kreuz{$\mathrm{sin}\,\beta = \frac{q}{p}$}

\dreieck{(0,0)}{(4,0)}{(0,3)}{p}{q}{r}{}{\beta}{}
\teil Rechter Winkel bei $C$. Welche Gleichung gilt? \hfill (P18) \\ \kreuz{$\mathrm{cos}\,\alpha = \frac{w}{v}$} \\ \kreuz{$\mathrm{cos}\,\alpha = \frac{v}{w}$} \\ \kreuz{$\mathrm{sin}\,\alpha = \frac{v}{w}$}

\dreieck{(0,0)}{(5,0)}{(3.2,2.4)}{u}{v}{w}{\alpha}{}{}
\teil Rechter Winkel bei $B$. Welche Gleichung gilt? \hfill (P18) \\ \kreuz{$\mathrm{tan}\,\gamma = \frac{m}{k}$} \\ \kreuz{$\mathrm{tan}\,\gamma = \frac{k}{m}$} \\ \kreuz{$\mathrm{sin}\,\gamma = \frac{m}{n}$}

\dreieck{(0,0)}{(5,0)}{(5,3)}{m}{n}{k}{}{}{\gamma}
\end{teile}
\end{aufgabe}

\begin{aufgabe}{Seite berechnen – Prüfungsaufgaben – weiter}
\begin{teile}
\teil Ein Grundstück $ABCD$ hat bei $A$ einen rechten Winkel. Die Diagonale $BD$ ist $840$ m lang, der Winkel $ADB$ beträgt $47^\circ$. Zeige rechnerisch, dass $AB \approx 614$ m ist. \hfill (P21)
\teil Ein Segelrevier $ABCD$ hat bei $A$ einen rechten Winkel. Die Diagonale $BD$ ist $2\,300$ m lang, der Winkel $ADB$ beträgt $41^\circ$. Zeige rechnerisch, dass $AB \approx 1\,509$ m ist. \hfill (P21)
\teil Ein Feld $ABCD$ hat bei $A$ einen rechten Winkel. Die Diagonale $BD$ ist $780$ m lang, der Winkel $ADB$ beträgt $53^\circ$. Berechne die Länge von $AD$. \hfill (P21) \\ AD ≈ \leerfeld[m]
\teil Ein Waldstück $ABCD$ hat bei $A$ einen rechten Winkel. Die Diagonale $BD$ ist $1\,650$ m lang, der Winkel $ADB$ beträgt $62^\circ$. Berechne die Länge von $AD$. \hfill (P21) \\ AD ≈ \leerfeld[m]
\teil Im Dachgeschoss ist eine senkrechte Wand $1{,}35$ m hoch. Die Dachschräge trifft den Boden im Punkt $P$ und bildet mit ihm einen Winkel von $38{,}5^\circ$. Wie weit ist $P$ von der Wand entfernt? \hfill (P17) \\ \leerfeld[m]
\teil In einem Zelt ist die senkrechte Seitenwand $0{,}85$ m hoch. Die Zeltbahn läuft bis zum Boden und bildet mit ihm einen Winkel von $27{,}4^\circ$. Wie weit liegt der Hering von der Seitenwand entfernt? \hfill (P17) \\ \leerfeld[m]
\teil Im Dreieck $ABC$ ist die Höhe von $C$ auf $AB$ eingezeichnet, sie ist $6{,}3$ m lang. Der Winkel bei $B$ beträgt $57^\circ$. Berechne die Seite $a = BC$. \hfill (P22) \\

\dreieck{(0,0)}{(6,0)}{(3.5,3)}{a}{b}{c}{}{57^\circ}{}
a ≈ \leerfeld[m]
\teil Im Dreieck $ABC$ ist die Höhe von $C$ auf $AB$ eingezeichnet, sie ist $4{,}8$ cm lang. Der Winkel bei $B$ beträgt $63^\circ$. Berechne die Seite $a = BC$. \hfill (P22) \\

\dreieck{(0,0)}{(5,0)}{(3.8,2.4)}{a}{b}{c}{}{63^\circ}{}
a ≈ \leerfeld[cm]
\teil $\mathrm{cos}\,60^\circ = \frac{4}{x}$ – berechne $x$. \hfill (P20) \\ x = \leerfeld
\teil $\mathrm{tan}\,45^\circ = \frac{6}{x}$ – berechne $x$. \hfill (P20) \\ x = \leerfeld
\teil Rechter Winkel bei $B$. Trage den Bruch ein. \hfill (P25) \\

\dreieck{(0,0)}{(5,0)}{(5,3)}{e}{f}{d}{}{}{\gamma}
$\mathrm{sin}\,\gamma =$ \leerfeld
\teil Rechter Winkel bei $A$. Trage den Bruch ein. \hfill (P25) \\

\dreieck{(1.8,2.4)}{(0,0)}{(5,0)}{z}{x}{y}{}{\beta}{}
$\mathrm{sin}\,\beta =$ \leerfeld
\teil Rechter Winkel bei $C$. Gib $\mathrm{tan}\,\alpha$ als Bruch an. \hfill (P20) \\

\dreieck{(4,0)}{(0,3)}{(0,0)}{q}{p}{s}{\alpha}{}{}
$\mathrm{tan}\,\alpha =$ \leerfeld
\teil Rechter Winkel bei $A$. Gib $\mathrm{tan}\,\gamma$ als Bruch an. \hfill (P20) \\

\dreieck{(0,0)}{(0,3)}{(5,0)}{w}{v}{u}{}{}{\gamma}
$\mathrm{tan}\,\gamma =$ \leerfeld
\end{teile}
\end{aufgabe}

\medskip
\uebersichtskasten{Seiten vom Winkel aus: Die Hypotenuse liegt dem rechten Winkel gegenüber, die Gegenkathete liegt dem markierten Winkel gegenüber, die Ankathete liegt an ihm – wechselt der Winkel, tauschen Gegen- und Ankathete. \\ Sinus, Kosinus, Tangens: sin α = Gegenkathete/Hypotenuse, cos α = Ankathete/Hypotenuse, tan α = Gegenkathete/Ankathete – der Taschenrechner (Grad-Modus) gibt den Wert zum Winkel. \\ Hypotenuse 12 cm, α = 40°: Gegenkathete = 12 · sin 40° ≈ 7,7 cm \quad Ankathete = 12 · cos 40° ≈ 9,2 cm \\ Steht die gesuchte Seite im Nenner, wird geteilt: Gegenkathete 9 cm, α = 40°: Hypotenuse = 9 : sin 40° ≈ 14,0 cm \quad Ankathete = 9 : tan 40° ≈ 10,7 cm \\ Merkregel: Kathete gesucht → Hypotenuse mal. Hypotenuse gesucht → Kathete geteilt. Erst am Ende runden.}
